% DERIVED from formal_layer.json — read-only, do not edit.
\begin{definitionv}{def:potential-path-space}[Potential-outcome space \((\mathcal S,\mathcal E)\)]
The public potential-outcome space \((\mathcal S,\mathcal E)\) is a measurable space of Borel functions from \([0,1]\) to \(\mathbb R\), satisfying the following conditions.
\begin{itemize}
\item \textbf{(Evaluation.)} The evaluation map
\[
e:\mathcal S\times[0,1]\longrightarrow\mathbb R,
\qquad e(f,a)=f(a),
\]
is jointly measurable for the product sigma algebra \(\mathcal E\otimes\mathcal B([0,1])\).
\item \textbf{(Extensionality.)} For \(f,g\in\mathcal S\),
\[
f=g
\quad\Longleftrightarrow\quad
e(f,a)=e(g,a)\ \text{for every }a\in[0,1].
\]
\item \textbf{(Threshold schedules.)} The space includes threshold schedules through the measurable map
\[
\iota:C([0,1],[0,1])\times[0,1]\longrightarrow\mathcal S,
\qquad
\iota(f,u)(a)=\mathbf1\{u\le f(a)\},
\]
where \(C([0,1],[0,1])\) carries its uniform-norm Borel sigma algebra.
\end{itemize}
The results apply to every fixed public space satisfying these conditions, with general schedules drawn from its Borel member functions. The potential-outcome schedule \(Y(\cdot)\) is a measurable random element of \(\mathcal S\), and its evaluation at dose \(a\) is
\[
Y(a)=e(Y(\cdot),a).
\]
\end{definitionv}

\begin{definitionv}{def:conditional-potential-mean}[Conditional potential mean \(\mu_x(a)\)]
The conditional potential-outcome mean is
\[
\mu_x(a)
=
E_P[Y(a)\mid X=x]
=
\frac{E_P[\mathbf1\{X=x\}Y(a)]}{p_x},
\qquad a\in[0,1],\quad x\in\mathcal X,\quad p_x>0.
\]
Here \(P\) is the structural probability law, \(X\) is the observed stratum coordinate, \(\mathcal X\) is the set of two observed strata, and
\[
p_x=P(X=x).
\]
The potential outcome \(Y(a)\) is the schedule evaluation in \cref{obj:def:potential-path-space}. Under the bounded-potential-outcome condition, the expectation exists and belongs to \([0,1]\). The model's stratum-overlap condition supplies \(p_x>0\) for both strata.
\end{definitionv}

\begin{assumptionv}{ass:stratum-overlap}[Uniform stratum overlap]
For a public constant \(p_{\min}\in(0,1/2]\), the stratum probabilities satisfy
\[
p_{\min}\le p_x\le 1-p_{\min}
\qquad\text{for every }x\in\mathcal X.
\]
\end{assumptionv}

\begin{assumptionv}{ass:latent-support}[Latent dose support]
The latent dose \(A\) satisfies
\[
A\in[0,1]\qquad\text{almost surely}.
\]
\end{assumptionv}

\begin{assumptionv}{ass:weak-design}[Local design density bounds]
For public constants \(0<c_{\mathrm{lo}}\le c_{\mathrm{hi}}\), each \(g_x\), \(x\in\mathcal X\), is a Lebesgue density satisfying
\[
c_{\mathrm{lo}}|t|^\kappa
\le g_x(t)\le
c_{\mathrm{hi}}|t|^\kappa
\qquad\text{for almost every }t\in[-1/2,1/2].
\]
\end{assumptionv}

\begin{assumptionv}{ass:holder-mean}[Hölder conditional mean increments]
For the potential-outcome schedules of \cref{obj:def:potential-path-space}, every stratum-specific conditional mean satisfies the radius-one Hölder increment condition
\[
\mu_x\in\mathcal H^\beta([0,1])
\quad\Longleftrightarrow\quad
|\mu_x(a)-\mu_x(a')|\le |a-a'|^\beta
\quad\text{for all }a,a'\in[0,1],
\qquad x\in\mathcal X.
\]
The class is defined by this increment condition; the unit-interval range is imposed separately.
\end{assumptionv}

\begin{assumptionv}{ass:mean-range}[Bounded conditional potential means]
For the potential-outcome schedules of \cref{obj:def:potential-path-space}, the conditional potential means satisfy
\[
0\le\mu_x(a)\le1
\qquad\text{for every }x\in\mathcal X
\text{ and }a\in[0,1].
\]
\end{assumptionv}

\begin{assumptionv}{ass:gaussian-channel}[Standard Gaussian error law]
The standardized error coordinate \(Z\) satisfies
\[
Z\sim N(0,1).
\]
\end{assumptionv}

\begin{assumptionv}{ass:error-independence}[Classical error independence]
The standardized error coordinate \(Z\) is independent of the observed stratum \(X\), latent dose \(A\), and potential-outcome schedule \(Y(\cdot)\):
\[
Z\perp(X,A,Y(\cdot)).
\]
\end{assumptionv}

\begin{assumptionv}{ass:conditional-unconfoundedness}[Conditional potential-outcome unconfoundedness]
The measurable random potential-outcome schedule \(Y(\cdot)\) and latent dose \(A\) are conditionally independent given the observed stratum \(X\):
\[
Y(\cdot)\perp A\mid X.
\]
\end{assumptionv}

\begin{assumptionv}{ass:bounded-potential-outcomes}[Bounded potential outcomes]
The potential outcomes in \cref{obj:def:potential-path-space} satisfy
\[
P\{Y(a)\in[0,1]\}=1
\qquad
\text{for every }a\in[0,1].
\]
\end{assumptionv}

\begin{assumptionv}{ass:consistency}[Realized outcome consistency]
The realized outcome \(Y\) equals the potential outcome from \cref{obj:def:potential-path-space} evaluated at the latent dose \(A\):
\[
Y=Y(A)\qquad P\text{-almost surely}.
\]
\end{assumptionv}

\begin{assumptionv}{ass:rank-uniform}[Conditional uniform rank law]
For every \(x\in\mathcal X\), the uniform-rank coordinate \(U\) satisfies
\[
U\mid X=x\sim\operatorname{Unif}[0,1].
\]
\end{assumptionv}

\begin{assumptionv}{ass:rank-treatment-independence}[Conditional rank–dose independence]
The rank coordinate \(U\) and latent dose \(A\) are conditionally independent given the observed stratum \(X\):
\[
U\perp A\mid X.
\]
\end{assumptionv}

\begin{assumptionv}{ass:structural-threshold}[Structural threshold schedules]
For the random potential-outcome schedule in \cref{obj:def:potential-path-space}, let \(U\) denote the uniform-rank coordinate and \(\mu_X(a)\) the conditional potential-outcome mean at the realized stratum. With \(P\)-probability one,
\[
Y(a)=\mathbf 1\{U\le \mu_X(a)\}
\qquad
\text{simultaneously for every }a\in[0,1].
\]
\end{assumptionv}

\begin{assumptionv}{ass:binary-potential-outcomes}[Binary potential outcomes]
The potential outcomes in \cref{obj:def:potential-path-space} satisfy
\[
P\{Y(a)\in\{0,1\}\}=1
\qquad
\text{for every }a\in[0,1].
\]
\end{assumptionv}

\begin{definitionv}{def:model-class}[Bounded-outcome structural model class]
The class \(\mathcal P_{K,\beta,\kappa,\sigma}\) consists of structural laws \(P\) of \((X,A,Z,Y(\cdot),Y,U)\), where \(X\) takes values in the two observed strata \(\mathcal X\), \(A\) is the latent dose, \(Z\) is the standardized error coordinate, \(Y(\cdot)\) is the random potential-outcome schedule on the public path space of \cref{obj:def:potential-path-space}, \(Y\) is the realized outcome, and \(U\) is an auxiliary rank coordinate. The parameters \(\beta,\kappa,\sigma\in\mathbb R\) are fixed, and the public class constants satisfy
\[
K=(p_{\min},c_{\mathrm{lo}},c_{\mathrm{hi}}),
\qquad
0<p_{\min}\le \tfrac12,
\qquad
0<c_{\mathrm{lo}}\le c_{\mathrm{hi}}.
\]
Membership in the class is defined by the following conditions:
\begin{itemize}
\item \textbf{(Stratum overlap.)} \cref{obj:ass:stratum-overlap}.
\item \textbf{(Latent support.)} \cref{obj:ass:latent-support}.
\item \textbf{(Conditional design.)} \cref{obj:ass:weak-design}, with envelope constants \(c_{\mathrm{lo}},c_{\mathrm{hi}}\) and exponent \(\kappa\).
\item \textbf{(Mean smoothness.)} \cref{obj:ass:holder-mean}, with exponent \(\beta\).
\item \textbf{(Mean range.)} \cref{obj:ass:mean-range}.
\item \textbf{(Gaussian channel.)} \cref{obj:ass:gaussian-channel}.
\item \textbf{(Error independence.)} \cref{obj:ass:error-independence}.
\item \textbf{(Conditional unconfoundedness.)} \cref{obj:ass:conditional-unconfoundedness}.
\item \textbf{(Bounded potential outcomes.)} \cref{obj:ass:bounded-potential-outcomes}.
\item \textbf{(Consistency.)} \cref{obj:ass:consistency}.
\end{itemize}
The contaminated dose \(W\) is defined by
\[
W=A+\sigma Z.
\]
The conditional potential-outcome mean functions \(\mu_x\) are defined in \cref{obj:def:conditional-potential-mean}. The lower-bound subclass additionally satisfies \cref{obj:ass:binary-potential-outcomes,obj:ass:rank-uniform,obj:ass:rank-treatment-independence,obj:ass:structural-threshold}.
\end{definitionv}

\begin{assumptionv}{ass:iid}[Independent and identically distributed observations]
The observed records \(O_1,\ldots,O_n\), each consisting of the observed stratum, contaminated dose, and realized outcome, are independent and identically distributed under the structural law \(P\).
\end{assumptionv}

\begin{definitionv}{def:observed-experiment}[Observed experiment \(Q_P^n\)]
The sample observed-law experiment is
\[
Q_P^n=\mathcal L_P(O_1,\ldots,O_n),
\]
where \(O_1,\ldots,O_n\) are independent and identically distributed observed records under the structural law \(P\), each consisting of the observed stratum, contaminated dose, and realized outcome.
\end{definitionv}

\begin{propositionv}{prop:identification}[Identification of the causal mean]
Fix the public potential-path space \((\mathcal S,\mathcal E)\) and evaluation map of \cref{obj:def:potential-path-space}, and public class constants \(K\) as in \cref{obj:def:model-class}. Let \(\mathcal X=\{0,1\}\) denote the two stratum labels, \(X\) the stratum coordinate, \(A\) the latent dose, \(Z\) the standardized error, and \(Y\) the realized outcome. Consider structural probability laws \(P\) and \(P'\) satisfying:
\begin{itemize}
\item \textbf{(Parameter ranges.)} \(\beta\in(0,1]\), \(\kappa\in[0,2]\), and \(\sigma\in[0,1/4]\).
\item \textbf{(Common model class.)} Both laws belong to the model class of \cref{obj:def:model-class} with the same \(K,\beta,\kappa,\sigma\) and public potential-path space.
\item \textbf{(Observed-law equality.)} Writing \(W=A+\sigma Z\) for the contaminated dose, the laws of the observed record coincide:
\[
\mathcal L_P(X,W,Y)=\mathcal L_{P'}(X,W,Y).
\]
\end{itemize}
For either structural law, write \(p_x(P)=P(X=x)\), and let \(\mu_x(a;P)\) denote its conditional potential-outcome mean from \cref{obj:def:conditional-potential-mean}. Set \(a_0=1/2\), the target dose, and \(t=A-a_0\), the centered latent dose. Define the population causal mean by
\[
\theta(P)=\int Y(a_0)\,dP,
\]
and analogously for \(P'\).

Then, for every \(x\in\mathcal X\), the stratum masses and conditional centered-dose measures agree:
\[
p_x(P)=p_x(P'),
\qquad
\mathcal L_P(A-a_0\mid X=x)
=
\mathcal L_{P'}(A-a_0\mid X=x).
\]
Their conditional latent mean measures also agree: for every Borel set \(B\subseteq\mathbb R\),
\[
\int_B \mu_x(t+a_0;P)\,
        \mathcal L_P(A-a_0\mid X=x)(dt)
=
\int_B \mu_x(t+a_0;P')\,
        \mathcal L_{P'}(A-a_0\mid X=x)(dt).
\]
Moreover,
\[
\mu_x(a;P)=\mu_x(a;P')
\qquad\text{for every }a\in[0,1].
\]
For every integer \(j\ge0\), the shrinking neighborhood has positive conditional latent mass:
\[
P\!\left(
-\frac1{j+3}<A-a_0<\frac1{j+3}
\,\middle|\,X=x
\right)>0,
\]
and its mean-measure to dose-measure ratio satisfies
\[
\lim_{j\to\infty}
\frac{
\displaystyle
\int_{(-1/(j+3),\,1/(j+3))}
\mu_x(t+a_0;P)\,
\mathcal L_P(A-a_0\mid X=x)(dt)
}{
\displaystyle
P\!\left(
-\frac1{j+3}<A-a_0<\frac1{j+3}
\,\middle|\,X=x
\right)
}
=
\mu_x(a_0;P).
\]
Finally,
\[
\theta(P)
=
\sum_{x\in\mathcal X}p_x(P)\mu_x(a_0;P),
\qquad
\theta(P)=\theta(P').
\]
\end{propositionv}

\begin{definitionv}{def:minimax-risk}[Minimax absolute-error risk $R_{K,n}(\sigma)$]
The minimax expected absolute error is
\[
R_{K,n}(\sigma)
=
\inf_T
\sup_{P\in\mathcal P_{K,\beta,\kappa,\sigma}}
E_{Q_P^n}\!\left[|T-\theta(P)|\right].
\]
Here \(\theta(P)\) is the population causal mean at the prespecified latent dose, and \(Q_P^n\) is the observed-sample law induced by the structural law \(P\), with potential-outcome schedules on the public path space of \cref{obj:def:potential-path-space}. The infimum ranges over measurable, possibly randomized estimators \(T\).
\end{definitionv}

\begin{definitionv}{def:honest-length}[Honest interval length $H_{K,\alpha,n}(\sigma)$]
The minimum worst-law expected length at fixed noncoverage probability \(\alpha\) is
\[
H_{K,\alpha,n}(\sigma)
=
\inf_{\substack{
C:\ \inf_{P\in\mathcal P_{K,\beta,\kappa,\sigma}}
Q_P^n\{\theta(P)\in C\}\ge 1-\alpha
}}
\sup_{P\in\mathcal P_{K,\beta,\kappa,\sigma}}
E_{Q_P^n}\!\left[\operatorname{length}(C)\right].
\]
The infimum ranges over measurable connected intervals \(C\). Here \(\theta(P)\) is the population causal mean at the prespecified latent dose, and \(Q_P^n\) is the observed-sample law induced by the structural law \(P\), with potential-outcome schedules on the public path space of \cref{obj:def:potential-path-space}.
\end{definitionv}

\begin{definitionv}{def:frontier-rate}[Frontier rate \(\rho_{n,\sigma}\)]
The frontier rate and its localization scale are
\[
\rho_{n,\sigma}=b_{n,\sigma}^{\beta},
\qquad
b_{n,\sigma}=
\begin{cases}
n^{-1/d},&0\le\sigma\le n^{-1/d},\\[2pt]
\dfrac{\sigma}{\sqrt{\log(e+n\sigma^d)}},
&n^{-1/d}<\sigma\le L_n^{-1/2},\\[6pt]
\dfrac{\log(e+\sigma^2L_n)}{L_n},
&L_n^{-1/2}<\sigma\le1/4.
\end{cases}
\]
Here the effective exponent \(d\) and logarithmic sample-size quantity \(L_n\) are
\[
d=2\beta+\kappa+1,
\qquad
L_n=\log(en).
\]
\end{definitionv}

\begin{theoremv}{thm:uniform-frontier}[Uniform risk and length frontier]*
Fix public class constants \(K=(p_{\min},c_{\mathrm{lo}},c_{\mathrm{hi}})\), smoothness exponent \(\beta\), design-decay exponent \(\kappa\), and noncoverage probability \(\alpha\), satisfying:
\begin{itemize}
\item \textbf{(Class constants.)} \(0<p_{\min}\le 1/2\) and \(0<c_{\mathrm{lo}}\le c_{\mathrm{hi}}\).
\item \textbf{(Parameter ranges.)} \(0<\beta\le1\), \(0\le\kappa\le2\), and \(0<\alpha<1/2\).
\item \textbf{(Envelope calibration.)} \(c_{\mathrm{lo}}\le(\kappa+1)2^\kappa\le c_{\mathrm{hi}}\).
\end{itemize}
Write \(R_n(\sigma)\) for the minimax absolute risk in \cref{obj:def:minimax-risk} and \(H_n(\sigma)\) for the minimum honest expected connected-interval length in \cref{obj:def:honest-length}, at these fixed parameters. Define the effective exponent \(d\), logarithmic sample scale \(L_n\), transition thresholds \(a_n,b_n\), and comparison scales \(D_n,F_n(\sigma),P_n(\sigma)\) by
\[
d=2\beta+\kappa+1,\qquad L_n=\log(en),\qquad
a_n=D_n=n^{-1/d},\qquad b_n=L_n^{-1/2},
\]
\[
F_n(\sigma)=\frac{\sigma}{\sqrt{\log(e+n\sigma^d)}},
\qquad
P_n(\sigma)=\frac{\log(e+\sigma^2L_n)}{L_n}.
\]
Define the localization scale \(b_{n,\sigma}\) and absolute-risk rate \(\rho_{n,\sigma}\) by
\[
b_{n,\sigma}=
\begin{cases}
D_n,&\sigma\le a_n,\\
F_n(\sigma),&a_n<\sigma\le b_n,\\
P_n(\sigma),&\sigma>a_n\ \text{and}\ \sigma>b_n,
\end{cases}
\qquad
\rho_{n,\sigma}=b_{n,\sigma}^{\beta}.
\]

There exist constants \(0<c<C\) and an integer \(n_0\), depending on \(K,\beta,\kappa,\alpha\), such that, for every public potential-path space as in \cref{obj:def:potential-path-space}, every integer \(n\ge n_0\), and every \(\sigma\in[0,1/4]\),
\[
c\rho_{n,\sigma}\le R_n(\sigma)\le C\rho_{n,\sigma},
\qquad
c\rho_{n,\sigma}\le H_n(\sigma)\le C\rho_{n,\sigma}.
\]
For the same \(C\) and every \(n\ge n_0\), the transition comparisons satisfy
\[
D_n\le C F_n(a_n),\qquad F_n(a_n)\le C D_n,
\]
\[
F_n(b_n)\le C P_n(b_n),\qquad P_n(b_n)\le C F_n(b_n).
\]

Moreover, for every real \(W\ge1\), there exist a constant \(C_K\ge1\) and an integer \(n_K\) such that, for every integer \(n\ge n_K\),
\[
0<a_n/W,\qquad Wa_n\le1/4,\qquad
0<b_n/W,\qquad Wb_n\le1/4,
\]
and the two transition windows satisfy
\[
C_K^{-1}D_n\le F_n(\sigma)\le C_KD_n
\qquad
\text{for every }\sigma\in[a_n/W,Wa_n],
\]
\[
C_K^{-1}P_n(\sigma)\le F_n(\sigma)\le C_KP_n(\sigma)
\qquad
\text{for every }\sigma\in[b_n/W,Wb_n].
\]
Here \(C_K\) and \(n_K\) may depend on \(W\) and the fixed parameters.

Finally, for every fixed \(\sigma\in(0,1/4]\), there exist constants \(0<c<C\) and an integer \(n_0\), which may depend on \(\sigma\) and the fixed parameters, such that, for every public potential-path space as in \cref{obj:def:potential-path-space} and every integer \(n\ge n_0\),
\[
c\left(\frac{\log\log n}{\log n}\right)^\beta
\le R_n(\sigma)\le
C\left(\frac{\log\log n}{\log n}\right)^\beta,
\]
\[
c\left(\frac{\log\log n}{\log n}\right)^\beta
\le H_n(\sigma)\le
C\left(\frac{\log\log n}{\log n}\right)^\beta.
\]
\end{theoremv}
% CAUSALSMITH-CITED-SCOPE-BEGIN thm:uniform-frontier
\verificationfootnotetext{\textbf{Formalization scope.} This result uses the cited conclusion \citep{PetrushevXu2005} (Theorem 2.4, equation (2.14)), \citep{https://people.math.sc.edu/pencho/Publications/kpx-06-28-06-web.pdf} (Theorem 2.2, equation (2.4)), \citep{https://dlmf.nist.gov/18.3} (Table 18.3.1, Jacobi row and its caption; orthogonality and squared norm.) and \citep{https://dlmf.nist.gov/5.11.E12} (Equation 5.11.12, restricted to the positive real axis.). The cited source proof is not formalized here: Lean verifies its use and all remaining steps. If that cited conclusion is false, the portions of this result that depend on it are not certified.}
% CAUSALSMITH-CITED-SCOPE-END thm:uniform-frontier

\begin{propositionv}{prop:error-free-reduction}[Error-free rates and benchmark comparison]*
Fix the public class constants \(K=(p_{\min},c_{\mathrm{lo}},c_{\mathrm{hi}})\) and parameters \(\beta,\kappa,\alpha\) satisfying:
\begin{itemize}
\item \textbf{(Class constants.)} \(0<p_{\min}\le 1/2\) and \(0<c_{\mathrm{lo}}\le c_{\mathrm{hi}}\).
\item \textbf{(Parameter ranges.)} \(0<\beta\le1\), \(0\le\kappa\le2\), and \(0<\alpha<1/2\), where \(\alpha\) is the noncoverage probability.
\item \textbf{(Envelope calibration.)}
\[
c_{\mathrm{lo}}\le(\kappa+1)2^\kappa\le c_{\mathrm{hi}}.
\]
\end{itemize}
There exist constants \(0<c<C\) and \(n_0\in\mathbb N\) such that, for every public potential-path space \((\mathcal S,\mathcal E)\) satisfying \cref{obj:def:potential-path-space} and every integer \(n\ge n_0\), the minimax absolute risk and minimum honest expected connected-interval length of \cref{obj:def:minimax-risk,obj:def:honest-length} satisfy
\[
c\,n^{-\beta/(2\beta+\kappa+1)}
\le R_n(0)\le
C\,n^{-\beta/(2\beta+\kappa+1)},
\qquad
c\,n^{-\beta/(2\beta+\kappa+1)}
\le H_n(0)\le
C\,n^{-\beta/(2\beta+\kappa+1)}.
\]

Moreover, fix any separate direct Gaussian-response benchmark with:
\begin{itemize}
\item \textbf{(Benchmark constants.)} \(c_{\mathrm G},r_{\mathrm G},M_{\mathrm G},\tau_{\mathrm G}>0\).
\item \textbf{(Evaluation neighborhood.)} \(x_*\in(0,1)\) and \(0<\delta\le\min\{x_*,1-x_*\}\).
\item \textbf{(Design.)} A fixed probability distribution \(\nu\) on \([0,1]\) with a nonnegative density \(g\) satisfying
\[
g(u)=c_{\mathrm G}|u-x_*|^\kappa
\quad\text{for }u\in[0,1]\text{ with }|u-x_*|\le\delta.
\]
\end{itemize}
The benchmark observations are \(B_i=f(D_i)+\xi_i\), \(i=1,\ldots,n\), where the \(D_i\) are independent draws from \(\nu\), and the \(\xi_i\) are independent \(N(0,\tau_{\mathrm G}^2)\) errors independent of the design. In the supremum below, \(f\) ranges over Borel functions on \([0,1]\) satisfying
\[
\|f\|_\infty\le M_{\mathrm G},
\qquad
\inf_{\substack{p\text{ a real polynomial}\\ \deg p\le\lfloor\beta\rfloor}}
\ \sup_{\substack{u\in[0,1]\\ |u-x_*|\le a}}
|f(u)-p(u-x_*)|
\le r_{\mathrm G}a^\beta
\]
for every
\[
0<a\le
\left\{\frac{\tau_{\mathrm G}^2(\kappa+1)}
{2c_{\mathrm G}r_{\mathrm G}^2n}\right\}^{1/(2\beta+\kappa+1)}.
\]
The infimum below ranges over measurable estimators from the benchmark observations, and \(E_f\) denotes expectation under their observed law. There exist constants \(0<c<C\) and \(n_0\in\mathbb N\), now also depending on the fixed benchmark, such that for every public potential-path space satisfying \cref{obj:def:potential-path-space} and every integer \(n\ge n_0\),
\[
c\,\inf_{\widehat f}\sup_f E_f|\widehat f(x_*)-f(x_*)|
\le R_n(0)\le
C\,\inf_{\widehat f}\sup_f E_f|\widehat f(x_*)-f(x_*)|.
\]
\end{propositionv}
% CAUSALSMITH-CITED-SCOPE-BEGIN prop:error-free-reduction
\verificationfootnotetext{\textbf{Formalization scope.} This result uses the cited conclusion \citep{PetrushevXu2005} (Theorem 2.4, equation (2.14)), \citep{https://people.math.sc.edu/pencho/Publications/kpx-06-28-06-web.pdf} (Theorem 2.2, equation (2.4)), \citep{https://dlmf.nist.gov/18.3} (Table 18.3.1, Jacobi row and its caption; orthogonality and squared norm.), \citep{https://dlmf.nist.gov/5.11.E12} (Equation 5.11.12, restricted to the positive real axis.) and \citep{Gaiffas2005} (Definition 2 and Theorem 1, equations (2.1), (2.4), and (2.5), PDF pages 3--4; exact power-law specialization with loss parameter one.). The cited source proof is not formalized here: Lean verifies its use and all remaining steps. If that cited conclusion is false, the portions of this result that depend on it are not certified.}
% CAUSALSMITH-CITED-SCOPE-END prop:error-free-reduction

\begin{definitionv}{synth_4}[Public expected-error score \(A_q(K)\)]
The public expected-error score of a weight and inverse-weight pair \((q,\ell_q)\) is
\[
A_q(K)
=
B_q(K)
+\frac{12}{b_q(K)}\sqrt{\frac{V_q(K)}{n}}
+2\sqrt{\frac3n}.
\]
Fix the public class constants \(K=(p_{\min},c_{\mathrm{lo}},c_{\mathrm{hi}})\), the mean-smoothness index \(\beta\), the design-decay index \(\kappa\), the Gaussian scale \(\sigma\), and the sample size \(n\). The pair satisfies the following conditions.
\begin{itemize}
\item \textbf{(Public constants.)}
\[
0<p_{\min}\le\frac12,
\qquad
0<c_{\mathrm{lo}}\le c_{\mathrm{hi}},
\qquad n\ge3.
\]
\item \textbf{(Inverse identity.)}
The weight \(q:[-1/2,1/2]\to[0,\infty)\) and its inverse \(\ell_q:\mathbb R\to\mathbb R\) are public Borel functions, and
\[
E_Z[\ell_q(t+\sigma Z)]=q(t),
\qquad -1/2\le t\le1/2,
\qquad Z\sim N(0,1),
\]
with every expectation absolutely finite.
\item \textbf{(Weighted integrability.)}
With weighted moments
\[
I_s(q)=\int_{-1/2}^{1/2}|t|^s q(t)\,dt,
\]
one has
\[
0<I_\kappa(q)<\infty,
\qquad
I_{\kappa+\beta}(q)<\infty.
\]
Moreover, \(\ell_q(V)\) is absolutely integrable under every law in \(\mathcal P_{K,\beta,\kappa,\sigma}\).
\item \textbf{(Second moment.)}
The public inverse-weight second-moment quantity is
\[
V_q(K)
=
c_{\mathrm{hi}}\int_{-1/2}^{1/2}
|t|^\kappa E_Z[\ell_q(t+\sigma Z)^2]\,dt
<\infty.
\]
The integral is a nonnegative extended integral.
\end{itemize}
The public denominator and localization-bias quantities are
\[
b_q(K)=p_{\min}c_{\mathrm{lo}}I_\kappa(q),
\qquad
B_q(K)=\frac{c_{\mathrm{hi}}}{c_{\mathrm{lo}}}
\frac{I_{\kappa+\beta}(q)}{I_\kappa(q)}.
\]
For a second-moment quantity expressed with factor \(16\), use the distinct notation
\[
V_q^{(16)}
=
16\int_{-1/2}^{1/2}|t|^\kappa
E_Z[\ell_q(t+\sigma Z)^2]\,dt
=
\frac{16}{c_{\mathrm{hi}}}V_q(K).
\]
Thus the second-moment term in the score has the equivalent expression
\[
\frac{12}{b_q(K)}\sqrt{\frac{V_q(K)}n}
=
\frac{12}{b_q(K)}
\sqrt{\frac{c_{\mathrm{hi}}V_q^{(16)}}{16n}}.
\]
Dependence on the fixed quantities \(\beta,\kappa,n,\sigma\) is suppressed in the notation.
\end{definitionv}

\begin{definitionv}{synth_2}[Polynomial localization weights and inverse-heat weights]
For a positive odd integer \(m\), define on \([-1/2,1/2]\)
\[
q_m^{\mathrm M}(t)=
\begin{cases}
\displaystyle
\left[\frac{\sin(m\arcsin(2t))}{2mt}\right]^6,
&t\ne0,\\[6pt]
1,&t=0.
\end{cases}
\]
For odd \(m\), this expression is a polynomial of degree \(6(m-1)\); use its polynomial extension to define \(q_m^{\mathrm M}\) on \(\mathbb R\). For known Gaussian error scale \(\sigma\in[0,1/4]\), define its companion observable inverse weight by
\[
\ell_m^{\mathrm M}(v)
=\sum_{j=0}^{3(m-1)}
\frac{(-\sigma^2/2)^j}{j!}
(q_m^{\mathrm M})^{(2j)}(v),
\qquad v\in\mathbb R,
\]
where the superscript \((2j)\) denotes the derivative of order \(2j\). For \(Z\sim N(0,1)\), this pair satisfies the exact Gaussian inversion identity
\[
E_Z\!\left[\ell_m^{\mathrm M}(t+\sigma Z)\right]
=q_m^{\mathrm M}(t),\qquad t\in\mathbb R.
\]
At sample size \(n\), the polynomial entries of the public dictionary use the odd indices
\[
m=2^j+1\le n^2,\qquad j\in\mathbb Z,\quad j\ge1.
\]
\end{definitionv}

\begin{definitionv}{synth_1}[Fourier localization weights and Gaussian inverses]
For \(0<h\le1/4\) and known Gaussian error scale \(\sigma\in[0,1/4]\), define
\[
Q(u)=
\begin{cases}
(\sin u/u)^6,&u\ne0,\\
1,&u=0,
\end{cases}
\qquad
q_h^{\mathrm F}(t)=Q(t/h),\quad t\in\mathbb R.
\]
Use the Fourier-transform convention
\[
\widehat f(\omega)=\int_{\mathbb R}e^{-i\omega t}f(t)\,dt,
\qquad
\mathcal F^{-1}[g](v)=\frac1{2\pi}
\int_{\mathbb R}e^{i\omega v}g(\omega)\,d\omega.
\]
The companion observable inverse weight is
\[
\ell_h^{\mathrm F}(v)
=\mathcal F^{-1}\!\left[
\widehat q_h^{\mathrm F}(\omega)e^{\sigma^2\omega^2/2}
\right](v),\qquad v\in\mathbb R.
\]
For \(Z\sim N(0,1)\), this pair satisfies the exact Gaussian inversion identity
\[
E_Z\!\left[\ell_h^{\mathrm F}(t+\sigma Z)\right]
=q_h^{\mathrm F}(t),\qquad t\in\mathbb R.
\]
At sample size \(n\), the Fourier entries of the public dictionary use the dyadic bandwidths
\[
h=2^{-j},\qquad j\in\mathbb Z,\quad j\ge2,
\qquad n^{-2}\le2^{-j}\le1/4.
\]
\end{definitionv}

\begin{definitionv}{def:weight-dictionary}[Inverse-weight dictionary \(\mathcal D_{n,\sigma}\)]
The finite public dictionary \(\mathcal D_{n,\sigma}\) contains the following weights, each paired with its exact inverse and assigned the public expected-error score \(A_q\):
\begin{itemize}
\item \textbf{(Fourier weights.)} The weights \(q_h^{\mathrm F}\), paired with \(\ell_h^{\mathrm F}\), for dyadic bandwidths \(h\in[n^{-2},1/4]\).
\item \textbf{(Polynomial weights.)} The weights \(q_m^{\mathrm M}\), paired with \(\ell_m^{\mathrm M}\), for odd indices \(m=2^j+1\le n^2\).
\item \textbf{(Constant weight.)} The constant weight \(1\), paired with the inverse \(1\).
\end{itemize}
\end{definitionv}

\begin{definitionv}{synth_6}[Public dictionary order]
The fixed dictionary order on \(\mathcal D_{n,\sigma}\) in \cref{obj:def:weight-dictionary} is the following order of tagged entries: the constant entry first, then the Fourier entries \(q_{2^{-j}}^{\mathrm F}\) in increasing integer index \(j\), then the polynomial entries \(q_{2^j+1}^{\mathrm M}\) in increasing integer index \(j\). The Fourier indices are those integers \(j\ge2\) satisfying \(n^{-2}\le2^{-j}\le1/4\); the polynomial indices are those integers \(j\ge1\) satisfying \(2^j+1\le n^2\). Each entry carries its paired inverse weight and public score \(A_q(K)\). Entries retain their tags even when their latent weight functions coincide. The first minimizer \(\widehat q\) is the earliest entry in this order attaining \(\min_{q\in\mathcal D_{n,\sigma}}A_q(K)\).
\end{definitionv}

\begin{definitionv}{synth_3}[Split-sample ratios]
Index the observed records \(O_i=(X_i,W_i,Y_i)\) by \(i=0,\ldots,n-1\), put \(V_i=W_i-a_0\) with target dose \(a_0=1/2\), and define
\[
\mathcal I_r=\{i\in\{0,\ldots,n-1\}:i\bmod3=r\},
\qquad r\in\{0,1,2\}.
\]
For a public latent weight \(q\) with observable inverse \(\ell_q\), let
\[
I_\kappa(q)=\int_{-1/2}^{1/2}|t|^\kappa q(t)\,dt,
\qquad
b_q(K)=p_{\min}c_{\mathrm{lo}}I_\kappa(q)>0,
\qquad K=(p_{\min},c_{\mathrm{lo}},c_{\mathrm{hi}}).
\]
When all three blocks are nonempty, define, for each \(x\in\mathcal X=\{0,1\}\),
\[
\widehat p_x
=\frac1{|\mathcal I_0|}\sum_{i\in\mathcal I_0}\mathbf1\{X_i=x\},
\]
\[
\widehat D_x(q)
=\frac1{|\mathcal I_1|}\sum_{i\in\mathcal I_1}\mathbf1\{X_i=x\}\ell_q(V_i),
\qquad
\widehat M_x(q)
=\frac1{|\mathcal I_2|}\sum_{i\in\mathcal I_2}\mathbf1\{X_i=x\}Y_i\ell_q(V_i).
\]
Floor the denominator at \(b_q(K)/2\) and project the ratio onto the public outcome range:
\[
\widehat\mu_x(q)
=\Pi_{[0,1]}\!\left(
\frac{\widehat M_x(q)}{\max\{\widehat D_x(q),b_q(K)/2\}}
\right),
\qquad
\Pi_{[0,1]}(u)=\min\{1,\max\{0,u\}\}.
\]
The corresponding total estimator is
\[
\widehat\theta_q
=
\begin{cases}
\displaystyle\sum_{x\in\mathcal X}\widehat p_x\widehat\mu_x(q),
&\min_{r\in\{0,1,2\}}|\mathcal I_r|>0,\\
1/2,&\min_{r\in\{0,1,2\}}|\mathcal I_r|=0.
\end{cases}
\]
This is the ratio construction used by \cref{obj:def:total-estimator}. The factor-\(16\) quantity used for the reference second-moment convention is kept distinct in \cref{obj:synth_4}; it does not impose \(p_{\min}c_{\mathrm{lo}}=1/16\) on the public class.
\end{definitionv}

\begin{algorithmv}{def:total-estimator}[Total estimator $\widehat\theta$]
The inputs are the observed records \(O_i=(X_i,W_i,Y_i)\), indexed by \(i=0,\ldots,n-1\), the public class constants \(K=(p_{\min},c_{\mathrm{lo}},c_{\mathrm{hi}})\), and the ordered dictionary \(\mathcal D_{n,\sigma}\) with scores \(A_q(K)\) from \cref{obj:synth_4}.
\begin{enumerate}
\item Form the sample blocks, their sizes, and the centered observed doses:
\[
\mathcal I_r
=\{i\in\{0,\ldots,n-1\}:i\bmod3=r\},
\qquad
N_r=|\mathcal I_r|,
\qquad r\in\{0,1,2\},
\]
\[
V_i=W_i-a_0,\qquad i=0,\ldots,n-1.
\]
\item If \(\min\{N_0,N_1,N_2\}=0\), output \(1/2\) and stop. The remaining steps apply when all three blocks are nonempty.
\item For each dictionary weight \(q\), use the public denominator bound
\[
I_\kappa(q)=\int_{-1/2}^{1/2}|t|^\kappa q(t)\,dt,
\qquad
b_q(K)=p_{\min}c_{\mathrm{lo}}I_\kappa(q).
\]
Compute, for each stratum \(x\in\mathcal X=\{0,1\}\),
\[
\widehat p_x
=\frac1{N_0}\sum_{i\in\mathcal I_0}\mathbf1\{X_i=x\},
\]
\[
\widehat D_x(q)
=\frac1{N_1}\sum_{i\in\mathcal I_1}
\mathbf1\{X_i=x\}\ell_q(V_i),
\qquad
\widehat M_x(q)
=\frac1{N_2}\sum_{i\in\mathcal I_2}
\mathbf1\{X_i=x\}Y_i\ell_q(V_i).
\]
\item Floor each denominator, project the ratio onto \([0,1]\) as in \cref{obj:synth_3}, and combine strata:
\[
\widehat\mu_x(q)
=\min\left\{1,\max\left\{0,
\frac{\widehat M_x(q)}
{\max\{\widehat D_x(q),b_q(K)/2\}}
\right\}\right\},
\qquad
\widehat\theta_q
=\sum_{x\in\mathcal X}\widehat p_x\widehat\mu_x(q).
\]
\item Select the first public-score minimizer in the fixed dictionary order:
\[
\widehat q
=\operatorname*{arg\,min}_{q\in\mathcal D_{n,\sigma}}A_q(K),
\]
with ties resolved by that order.
\item Output \(\widehat\theta=\widehat\theta_{\widehat q}\).
\end{enumerate}
\end{algorithmv}

\begin{theoremv}{thm:observable-upper}[Observable estimator and comparison bounds]
Fix the following public parameters:
\begin{itemize}
\item \textbf{(Class constants.)} \(K\) specifies a minimum stratum mass in \((0,1/2]\) and positive lower and upper weak-design envelope constants, with the lower constant at most the upper constant.
\item \textbf{(Smoothness.)} The smoothness exponent satisfies \(\beta\in(0,1]\).
\item \textbf{(Design decay.)} The design-decay exponent satisfies \(\kappa\in[0,2]\).
\end{itemize}
There exist \(C>0\) and \(n_0\in\mathbb N\), depending on \(K,\beta,\kappa\), such that the following conclusions hold uniformly over every measurable schedule space and potential-path structure in \cref{obj:def:potential-path-space}, every \(n\ge n_0\), and every public error scale \(\sigma\in[0,1/4]\). Let \(L_n=\log(en)\) be the logarithmic sample scale and let \(\rho_{n,\sigma}\) be the absolute-risk frontier in \cref{obj:def:frontier-rate}.

For every structural probability law \(P\) in the model class of \cref{obj:def:model-class} with parameters \(K,\beta,\kappa,\sigma\), under the sampling condition in \cref{obj:ass:iid}, the observed experiment \(Q_P^n\) in \cref{obj:def:observed-experiment} and the selected estimator \(\widehat\theta\) in \cref{obj:def:total-estimator} satisfy
\[
\mathbb E_{Q_P^n}\!\left[
  \left|\widehat\theta-\theta(P)\right|
\right]
\le C\rho_{n,\sigma},
\qquad
\theta(P)=\mathbb E_P[Y(a_0)].
\]
Here \(\theta(P)\) is the population causal mean, \(a_0\) is the prespecified evaluation dose, and \(Y(a_0)\) is the potential outcome obtained by evaluating the random schedule there.

The public dictionary \(\mathcal D_{n,\sigma}\) in \cref{obj:def:weight-dictionary} contains comparison witnesses as follows:
\begin{itemize}
\item \textbf{(Fourier range.)} If \(\sigma\le L_n^{-1/2}\), there is a Fourier entry indexed by a nonnegative integer whose associated inverse pair \((q,\ell_q)\) has public expected-error score \(A_q\) satisfying
\[
A_q\le C\rho_{n,\sigma}.
\]
\item \textbf{(Polynomial range.)} If \(L_n^{-1/2}<\sigma\), there is a polynomial entry indexed by a nonnegative integer whose associated inverse pair \((q,\ell_q)\) has public expected-error score \(A_q\) satisfying
\[
A_q\le C\rho_{n,\sigma}.
\]
\end{itemize}
The pairs and scores are those associated with the entries in \cref{obj:def:weight-dictionary}. For each witness, under the usual measurability and integrability conditions, the latent weight \(q\) and its Gaussian inverse weight \(\ell_q\) satisfy
\[
q(t)\ge0
\quad\text{for every }t\in[-1/2,1/2],
\qquad
0<I_\kappa(q)
:=\int_{-1/2}^{1/2}|t|^\kappa q(t)\,dt<\infty,
\qquad
V_q<\infty,
\]
where \(V_q\) denotes the separately normalized factor-\(16\) second-moment certificate of \cref{obj:lem:public-error-certificate}. Their Gaussian inversion identity holds pointwise:
\[
\int_{\mathbb R}\ell_q(t+\sigma z)
       \frac{e^{-z^2/2}}{\sqrt{2\pi}}\,dz
=q(t)
\quad\text{for every }t\in[-1/2,1/2].
\]
\end{theoremv}

\begin{definitionv}{def:honest-interval}[Score-based interval \(C_{n,\sigma}\)]
The connected interval is
\[
C_{n,\sigma}
=
\begin{cases}
\bigl[\widehat\theta-A_{\widehat q}(K)/\alpha,\,
      \widehat\theta+A_{\widehat q}(K)/\alpha\bigr]\cap[0,1],
      &\min_{r\in\{0,1,2\}}|\mathcal I_r|>0,\\
[0,1],&\min_{r\in\{0,1,2\}}|\mathcal I_r|=0.
\end{cases}
\]
Here \(K\) is public, \(\alpha>0\) is the noncoverage level, and \(\widehat\theta\) is the selected-weight estimator using the blocks
\[
\mathcal I_r
=\{i\in\{0,\ldots,n-1\}:i\bmod 3=r\},
\qquad r\in\{0,1,2\},
\]
and the first minimizer \(\widehat q\) of \(A_q(K)\) in the fixed dictionary order.
\end{definitionv}

\begin{theoremv}{thm:finite-honesty}[Finite-sample honesty and length]
Fix public constants \(K=(p_{\min},c_{\mathrm{lo}},c_{\mathrm{hi}})\) and parameters satisfying:
\begin{itemize}
\item \textbf{(Class constants.)} \(0<p_{\min}\le 1/2\) and \(0<c_{\mathrm{lo}}\le c_{\mathrm{hi}}\).
\item \textbf{(Smoothness.)} \(0<\beta\le 1\).
\item \textbf{(Design decay.)} \(0\le\kappa\le 2\).
\item \textbf{(Honesty level.)} \(0<\alpha<1\), where \(\alpha\) is the noncoverage probability.
\end{itemize}
For any public potential-outcome path space as in \cref{obj:def:potential-path-space}, let \(P\) range over the structural laws in \(\mathcal P_{K,\beta,\kappa,\sigma}\) of \cref{obj:def:model-class}. Write
\[
\theta(P)=\mathbb E_P[Y(a_0)],
\qquad
\rho_{n,\sigma}=b_{n,\sigma}^{\beta},
\]
where \(a_0\) is the prespecified latent evaluation dose and \(b_{n,\sigma}\) is the localization scale in \cref{obj:def:frontier-rate}. Let \(Q_P^n\) be the observed sample law in \cref{obj:def:observed-experiment}, and let \(C_{n,\sigma}\) be the interval in \cref{obj:def:honest-interval}, constructed with \(K,\beta,\kappa,\alpha\).

For every such path space, every sample size \(n\ge 0\), every \(\sigma\in[0,1/4]\), and every \(P\in\mathcal P_{K,\beta,\kappa,\sigma}\),
\[
Q_P^n\!\left\{\theta(P)\in C_{n,\sigma}\right\}\ge 1-\alpha.
\]
Moreover, there exist \(C=C(K,\beta,\kappa,\alpha)>0\) and an integer \(n_0=n_0(K,\beta,\kappa,\alpha)\ge 0\) such that, simultaneously for every such path space, every \(n\ge n_0\), every \(\sigma\in[0,1/4]\), and every \(P\in\mathcal P_{K,\beta,\kappa,\sigma}\),
\[
\mathbb E_{Q_P^n}\!\left[\operatorname{length}(C_{n,\sigma})\right]
\le C\,\rho_{n,\sigma}.
\]
\end{theoremv}

\begin{lemmav}{lem:realized-dose-mean}[Realized dose mean identities]
Suppose the following conditions hold:
\begin{itemize}
\item \textbf{(Public model primitives.)} Fix a potential-outcome path space as in \cref{obj:def:potential-path-space} and public class constants \(K\) as in \cref{obj:def:model-class}.
\item \textbf{(Parameter ranges.)} Let \(\beta\in(0,1]\), \(\kappa\in[0,2]\), and \(\sigma\in[0,1/4]\).
\item \textbf{(Structural law.)} Let \(P\) be a structural probability law in the model class \(\mathcal P_{K,\beta,\kappa,\sigma}\) of \cref{obj:def:model-class}. Write \(\mathcal X=\{0,1\}\) for the stratum space, and let \(X\in\mathcal X\), \(A\in\mathbb R\), \(Z\in\mathbb R\), and \(Y\in\mathbb R\) denote its stratum, latent dose, classical error, and realized outcome coordinates, respectively.
\end{itemize}
With \(\mu_x(a)\) denoting the conditional potential-outcome mean defined in \cref{obj:def:conditional-potential-mean}, the following statements hold \(P\)-almost surely:
\[
0\le Y\le1,\qquad
E_P[Y\mid A,X,Z]=\mu_X(A),\qquad
E_P[Y\mid A,X]=\mu_X(A).
\]
Moreover, for every bounded measurable function
\(F:\mathbb R\times\mathcal X\times\mathbb R\to\mathbb R\),
\[
E_P\!\left[F(A,X,Z)Y\right]
=
E_P\!\left[F(A,X,Z)\mu_X(A)\right].
\]
\end{lemmav}

\begin{lemmav}{lem:compact-gaussian-convolution-injective}[Injectivity of compact Gaussian convolution]
Let \(\sigma\in\mathbb R\), and let \(\Phi_\sigma\) denote the centered Gaussian probability measure with variance \(\sigma^2\), with \(\Phi_0=\delta_0\), the unit point mass at zero. Let \(\mu_1,\mu_2,\nu_1,\nu_2\) be nonnegative Borel measures on \(\mathbb R\). Suppose that:
\begin{itemize}
\item \textbf{(Error scale.)} \(0\le \sigma\le 1/4\).
\item \textbf{(Finiteness and support.)} Each of \(\mu_1,\mu_2,\nu_1,\nu_2\) is finite and assigns zero mass to \(\mathbb R\setminus[-1/2,1/2]\).
\item \textbf{(Convolution equality.)}
\[
\mu_1*\Phi_\sigma+\nu_2*\Phi_\sigma
=
\nu_1*\Phi_\sigma+\mu_2*\Phi_\sigma.
\]
\end{itemize}
Then
\[
\mu_1+\nu_2=\nu_1+\mu_2.
\]
\end{lemmav}

\begin{lemmav}{lem:positive-ratio}[Positive denominators and approximation]
Fix a potential-path space \((\mathcal S,\mathcal E)\) as in \cref{obj:def:potential-path-space}. Let \(\mathcal X\) consist of two strata, and let \(P\) be a structural probability law with coordinates \(X\in\mathcal X\), \(A,Z,Y,U\in\mathbb R\), and potential-outcome schedule \(Y(\cdot)\in\mathcal S\). Set
\[
a_0=\frac12,\qquad t=A-a_0,\qquad W=A+\sigma Z,\qquad V=W-a_0=t+\sigma Z.
\]
Let \(q:\mathbb R\to\mathbb R\) be a public latent weight and \(\ell_q:\mathbb R\to\mathbb R\) its observable inverse weight. Define the weighted moments by
\[
I_\gamma(q)=\int_{-1/2}^{1/2}|t|^\gamma q(t)\,dt.
\]
Suppose the following conditions hold:
\begin{itemize}
\item \textbf{(Public constants.)} The constants \(K=(p_{\min},c_{\mathrm{lo}},c_{\mathrm{hi}})\) satisfy
\[
0<p_{\min}\le\frac12,\qquad 0<c_{\mathrm{lo}}\le c_{\mathrm{hi}}.
\]
\item \textbf{(Parameter ranges.)}
\[
\beta\in(0,1],\qquad \kappa\in[0,2],\qquad \sigma\in[0,1/4].
\]
\item \textbf{(Structural model.)} The law \(P\) satisfies \cref{obj:def:model-class} with path space \((\mathcal S,\mathcal E)\), constants \(K\), and parameters \((\beta,\kappa,\sigma)\).
\item \textbf{(Admissible weight.)} The functions \(q\) and \(\ell_q\) are measurable, and
\[
q(t)\ge0\quad\text{for every }t\in[-1/2,1/2],
\qquad 0<I_\kappa(q)<\infty.
\]
\item \textbf{(Exact Gaussian inversion.)} For every \(t\in[-1/2,1/2]\), the function \(z\mapsto\ell_q(t+\sigma z)\) is absolutely integrable under the standard Gaussian law, and
\[
\int_{\mathbb R}\ell_q(t+\sigma z)\frac{e^{-z^2/2}}{\sqrt{2\pi}}\,dz=q(t).
\]
The inverse weight is also absolutely integrable under the structural law:
\[
E_P[|\ell_q(V)|]<\infty.
\]
\end{itemize}
For every \(x\in\mathcal X\), the denominator is strictly positive and satisfies
\[
0<E_P[\mathbf1\{X=x\}\ell_q(V)],
\qquad
E_P[\mathbf1\{X=x\}\ell_q(V)]
\ge p_{\min}c_{\mathrm{lo}}I_\kappa(q).
\]
Consequently, the localized population ratio
\[
m_{x,q}
=
\frac{E_P[\mathbf1\{X=x\}Y\ell_q(V)]}
     {E_P[\mathbf1\{X=x\}\ell_q(V)]}
\]
is well defined and, with \(\mu_x(a_0)\) as in \cref{obj:def:conditional-potential-mean}, obeys
\[
|m_{x,q}-\mu_x(a_0)|
\le
\frac{c_{\mathrm{hi}}}{c_{\mathrm{lo}}}
\frac{I_{\kappa+\beta}(q)}{I_\kappa(q)}.
\]
These conclusions hold throughout the stated parameter ranges, including \(\sigma=0\).
\end{lemmav}

\begin{lemmav}{lem:public-error-certificate}[Public expected-error certificate]
Fix the public potential-path space of \cref{obj:def:potential-path-space} and public class constants \(K\) satisfying the restrictions in \cref{obj:def:model-class}. Suppose that:
\begin{itemize}
\item \textbf{(Parameter ranges.)} \(\beta\in(0,1]\), \(\kappa\in[0,2]\), \(n\ge3\), and \(\sigma\in[0,1/4]\).
\item \textbf{(Structural law.)} The structural law \(P\) belongs to \(\mathcal P_{K,\beta,\kappa,\sigma}\), the model class with constants \(K\) defined in \cref{obj:def:model-class}.
\item \textbf{(Sampling.)} The observed sample satisfies \cref{obj:ass:iid}, with joint law \(Q_P^n\) defined in \cref{obj:def:observed-experiment}.
\item \textbf{(Inverse pair.)} Let \(q\) be a public latent weight and \(\ell_q\) its observable inverse weight. Under the usual measurability and integrability conditions, they satisfy
\[
q(t)\ge0\qquad(t\in[-1/2,1/2]),
\qquad
0<I_\kappa(q):=\int_{-1/2}^{1/2}|t|^\kappa q(t)\,dt<\infty,
\]
and the pointwise Gaussian inversion identity
\[
\int_{\mathbb R}\ell_q(t+\sigma z)
       \frac{e^{-z^2/2}}{\sqrt{2\pi}}\,dz
=q(t)
\qquad(t\in[-1/2,1/2]),
\]
where the Gaussian integrand is absolutely integrable for every such \(t\).
\item \textbf{(Law integrability.)} Let \(W\) denote the contaminated dose and \(a_0\) the prespecified evaluation dose, and define the observed-dose displacement by
\[
V:=W-a_0.
\]
Then
\[
\mathbb E_P[|\ell_q(V)|]<\infty.
\]
\item \textbf{(Public second moment.)} Define the public inverse-weight second moment, as an extended nonnegative integral, by
\[
V_q:=
16\int_{-1/2}^{1/2}|t|^\kappa
   \left(
   \int_{\mathbb R}\ell_q(t+\sigma z)^2
       \frac{e^{-z^2/2}}{\sqrt{2\pi}}\,dz
   \right)dt.
\]
Assume \(V_q<\infty\).
\end{itemize}
Let \(\widehat\theta_q\) be the weight-indexed estimator of \cref{obj:def:total-estimator}, and let \(A_q(K)\) be its public expected-error score from that construction, evaluated at the fixed parameters \(\beta,\kappa,n,\sigma\). Define the population causal target by
\[
\theta(P):=\mathbb E_P[Y(a_0)],
\]
where \(Y(a_0)\) is the potential outcome at the evaluation dose in \cref{obj:def:potential-path-space}. Then
\[
\mathbb E_{Q_P^n}
 \left[|\widehat\theta_q-\theta(P)|\right]
\le A_q(K).
\]
\end{lemmav}

\begin{lemmav}{lem:dictionary-score-rate}[Dictionary admissibility and score rate]
Fix the following public objects and parameters:
\begin{itemize}
\item \textbf{(Path space.)} A potential-path space \((\mathcal S,\mathcal E)\) as in \cref{obj:def:potential-path-space}.
\item \textbf{(Class constants.)} \(K=(p_{\min},c_{\mathrm{lo}},c_{\mathrm{hi}})\), with
\[
0<p_{\min}\le \tfrac12,
\qquad
0<c_{\mathrm{lo}}\le c_{\mathrm{hi}}.
\]
\item \textbf{(Exponents.)} \(\beta\in(0,1]\) and \(\kappa\in[0,2]\).
\end{itemize}
Use the public dictionary \(\mathcal D_{n,\sigma}\) from \cref{obj:def:weight-dictionary} and the public expected-error scores \(A_q(K)\) from \cref{obj:synth_4}. For every integer \(n\ge1\), every \(\sigma\in[0,1/4]\), and every inverse pair \((q,\ell_q)\in\mathcal D_{n,\sigma}\), both weights are measurable, \(q\) is nonnegative on \([-1/2,1/2]\), and its design-weighted moment satisfies
\[
0<I_\kappa(q)
=
\int_{-1/2}^{1/2}|t|^\kappa q(t)\,dt
<\infty.
\]
For every \(t\in[-1/2,1/2]\), the exact Gaussian inversion identity holds with absolute integrability:
\[
\int_{\mathbb R}|\ell_q(t+\sigma z)|\,\mathcal N(0,1)(dz)<\infty,
\qquad
\int_{\mathbb R}\ell_q(t+\sigma z)\,\mathcal N(0,1)(dz)=q(t),
\]
where \(\mathcal N(0,1)\) denotes the standard Gaussian probability measure. The inverse-weight second-moment integral satisfies
\[
16\int_{-1/2}^{1/2}|t|^\kappa
\left(\int_{\mathbb R}\ell_q(t+\sigma z)^2\,\mathcal N(0,1)(dz)\right)\,dt
<\infty.
\]
Let \(V\) denote the observed-dose displacement from the prespecified evaluation dose \(a_0\), and let \(W\) denote the contaminated dose:
\[
V=W-a_0,
\qquad
W=A+\sigma Z,
\]
where \(A\) is the latent dose and \(Z\) is the standardized Gaussian error. For every structural law \(P\in\mathcal P_{K,\beta,\kappa,\sigma}\) from \cref{obj:def:model-class}, the inverse weight is absolutely integrable under the observed centered dose:
\[
\mathbb E_P[|\ell_q(V)|]<\infty.
\]

Let \(\widehat q\) be the first public-score minimizer in the fixed dictionary order from \cref{obj:def:total-estimator}, and let \(\rho_{n,\sigma}\) and \(L_n\) be the frontier rate and logarithmic sample-size quantity from \cref{obj:def:frontier-rate}. There exist \(C>0\) and \(n_0\in\mathbb N\), uniform in \(\sigma\in[0,1/4]\), such that for every \(n\ge n_0\) and every \(\sigma\in[0,1/4]\),
\[
A_{\widehat q}(K)\le C\rho_{n,\sigma}.
\]
For the same \(C\) and \(n_0\), whenever \(\sigma\le L_n^{-1/2}\), there exists \(k\in\mathbb N\) such that the Fourier inverse pair indexed by \(k\) in \cref{obj:def:weight-dictionary} belongs to \(\mathcal D_{n,\sigma}\) and its public score is at most \(C\rho_{n,\sigma}\). Whenever \(L_n^{-1/2}<\sigma\), there exists \(j\in\mathbb N\) such that the polynomial inverse pair with index
\[
m=2^j+1
\]
belongs to \(\mathcal D_{n,\sigma}\) and satisfies
\[
A_{q_m^{\mathrm M}}(K)\le C\rho_{n,\sigma}.
\]
\end{lemmav}

\begin{definitionv}{def:packet-handle}[Filtered Jacobi packet \(\psi_m\)]
For each integer \(m\ge4\), the zero-extended filtered Jacobi packet is
\[
\psi_m(u)=
\begin{cases}
(1-u^2)^r K_m(2u^2-1)/K_m(-1),
& |u|\le1,\ \kappa>0,\\
(1-u^2)^r\widetilde K_m(u)/\widetilde K_m(0),
& |u|\le1,\ \kappa=0,\\
0,& |u|>1.
\end{cases}
\]
The fixed taper parameter and smooth spectral filter are
\[
r=4,\qquad
\eta(s)=
\begin{cases}
\exp\{-1/((s-9/8)(15/8-s))\},
& 9/8<s<15/8,\\
0,& \text{otherwise}.
\end{cases}
\]
The filter is nonnegative, supported inside \((1,2)\), and positive on \([5/4,7/4]\). Write \(\mathsf J_j^{(a,b)}\) for the standard Jacobi polynomial of degree \(j\), and define its squared norm by
\[
H_j^{a,b}
=
\int_{-1}^{1}
\bigl(\mathsf J_j^{(a,b)}(z)\bigr)^2
(1-z)^a(1+z)^b\,dz.
\]
For \(\kappa>0\), the endpoint parameter and filtered sum are
\[
b=\frac{\kappa-1}{2},\qquad
K_m(z)=\sum_{j\ge0}\eta(j/m)
\frac{\mathsf J_j^{(r,b)}(-1)\mathsf J_j^{(r,b)}(z)}
{H_j^{r,b}}.
\]
For \(\kappa=0\), the symmetric interior sum is
\[
\widetilde K_m(u)=\sum_{j\ge0}\eta(j/m)
\frac{\mathsf J_j^{(r,r)}(0)\mathsf J_j^{(r,r)}(u)}
{H_j^{r,r}}.
\]
\end{definitionv}

\begin{lemmav}{lem:jacobi-binomial}[Jacobi polynomial binomial representation]
Let \(a,b\in\mathbb R\) satisfy \(a,b>-1/2\), let \(j\ge0\) be an integer, and let \(z\in\mathbb R\). Write \(P_j^{(a,b)}=\mathsf J_j^{(a,b)}\) for the standard Jacobi polynomial. Define the generalized binomial coefficient, for \(v\in\mathbb R\) and integers \(i\ge0\), by
\[
\binom{v}{i}=\frac{\prod_{k=0}^{i-1}(v-k)}{i!},
\qquad \binom{v}{0}=1,
\]
where the empty product equals \(1\). Then
\[
P_j^{(a,b)}(z)
=2^{-j}\sum_{i=0}^{j}
\binom{j+a}{i}\binom{j+b}{j-i}
(z-1)^{j-i}(z+1)^i.
\]
\end{lemmav}

\begin{theoremv}{thm:weighted-packet-construction}[Weighted packet construction]*
For every fixed \(\kappa\in[0,2]\), the degeneracy exponent, there exist real constants \(C>0\) and \(c>0\) such that, for every integer \(m\ge4\), the filtered-Jacobi construction has positive normalizing value:
\[
\begin{cases}
K_m(-1)>0, & \kappa>0,\\
\widetilde K_m(0)>0, & \kappa=0.
\end{cases}
\]
Here \(K_m(-1)\) and \(\widetilde K_m(0)\) are respectively the endpoint and symmetric interior normalizing values of the construction. The resulting normalized, tapered, zero-extended packet \(\psi_m\) is even and continuously differentiable on \(\mathbb R\), and satisfies
\[
\psi_m(u)=0\quad\text{for }u\notin[-1,1],
\qquad
\psi_m(0)=1,
\]
\[
|\psi_m(u)|\le C,\qquad
|\psi_m'(u)|\le Cm
\quad\text{for every }u\in\mathbb R,
\]
\[
\int_{-1}^{1}|u|^\kappa|\psi_m(u)|\,du
\le Cm^{-\kappa-1},
\qquad
\int_{-1}^{1}|u|^\kappa\psi_m(u)^2\,du
\le Cm^{-\kappa-1},
\]
and
\[
\int_{-1}^{1}|u|^\kappa u^j\psi_m(u)\,du=0
\quad\text{for every integer }0\le j<cm.
\]
\end{theoremv}
% CAUSALSMITH-CITED-SCOPE-BEGIN thm:weighted-packet-construction
\verificationfootnotetext{\textbf{Formalization scope.} This result uses the cited conclusion \citep{PetrushevXu2005} (Theorem 2.4, equation (2.14)), \citep{https://people.math.sc.edu/pencho/Publications/kpx-06-28-06-web.pdf} (Theorem 2.2, equation (2.4)), \citep{https://dlmf.nist.gov/18.3} (Table 18.3.1, Jacobi row and its caption; orthogonality and squared norm.) and \citep{https://dlmf.nist.gov/5.11.E12} (Equation 5.11.12, restricted to the positive real axis.). The cited source proof is not formalized here: Lean verifies its use and all remaining steps. If that cited conclusion is false, the portions of this result that depend on it are not certified.}
% CAUSALSMITH-CITED-SCOPE-END thm:weighted-packet-construction

\begin{definitionv}{synth_5}[Structural lower-bound witness laws]
Fix \(0<\beta\le1\), \(0\le\kappa\le2\), \(0\le\sigma\le1/4\), a positive sample size \(n\), and the public measurable schedule space \((\mathcal S,\mathcal E)\) with threshold embedding \(\iota\). Put \(a_0=1/2\) and \(d=2\beta+\kappa+1\). For public class constants \(K=(p_{\min},c_{\mathrm{lo}},c_{\mathrm{hi}})\), require
\[
0<p_{\min}\le1/2,\qquad
c_{\mathrm{lo}}\le(\kappa+1)2^\kappa\le c_{\mathrm{hi}}.
\]
The three witness laws share independent seed coordinates \(X,T,Z,U\), where \(P\{X=0\}=P\{X=1\}=1/2\), \(Z\sim N(0,1)\), \(U\sim\operatorname{Unif}[0,1]\), and
\[
g_\kappa(t)=(\kappa+1)2^\kappa|t|^\kappa\mathbf1\{|t|\le1/2\}
\]
is the density of \(T\). Set \(A=a_0+T\).

Let \(\epsilon,c_p,C>0\). The tuning data consist of \(h,\ell\in(0,1/4]\) and nonnegative integers \(m,J\). In the direct regime \(\sigma\le n^{-1/d}\), set
\[
h_0=n^{-1/d},\qquad h=\ell=h_0,\qquad m=J=0,
\]
and \(\delta(t)=\epsilon h_0^\beta(1-(t/h_0)^2)^2\mathbf1\{|t|\le h_0\}\). In either inverse regime \(\sigma>n^{-1/d}\), use the normalized filtered-Jacobi packet \(\psi_m\) and require
\[
m\ge4,\qquad h=\ell/m,\qquad J=\lfloor c_pm\rfloor\ge1,
\qquad \delta(t)=\epsilon h^\beta\psi_m(t/\ell).
\]
The perturbation is supported on \([-\ell,\ell]\), is integrable against \(g_\kappa\), and satisfies
\[
\int_{\mathbb R}t^jg_\kappa(t)\delta(t)\,dt=0\qquad(0\le j<J).
\]
The continuous profiles
\[
\mu_\star(a)=1/2,\qquad \mu_\pm(a)=1/2\pm\delta(a-a_0)
\]
take values in \([1/8,7/8]\) and obey \(|\mu_\pm(a)-\mu_\pm(a')|\le|a-a'|^\beta\).

For \(s\in\{+,-,\star\}\), define \(P_s\) as the six-coordinate structural law of
\[
\bigl(X,A,Z,\iota(\mu_s,U),\mathbf1\{U\le\mu_s(A)\},U\bigr).
\]
Thus \(Y_s(a)=\mathbf1\{U\le\mu_s(a)\}\) simultaneously for all \(a\), \(Y_s=Y_s(A)\), and the observed projection is \((X,W,Y_s)\) with \(W=A+\sigma Z\). The retained rank \(U\) gives the common-seed representation required by the structural laws. The three laws have the same \((X,A)\) marginal and satisfy
\[
\mu_++\mu_-=2\mu_\star,\qquad
\theta(P_\pm)=1/2\pm\delta(0),\qquad \theta(P_\star)=1/2.
\]

Let \(Q_s=\mathcal L_{P_s}(X,W,Y_s)\). For every Borel set \(B\subseteq\mathbb R\) and \(x\in\{0,1\}\),
\[
Q_s\{X=x,W-a_0\in B,Y_s=1\}
=\frac12\int_{-1/2}^{1/2}g_\kappa(t)\mu_s(a_0+t)\Phi_\sigma(B-t)\,dt,
\]
with the analogous formula containing \(1-\mu_s\) when \(Y_s=0\). Here \(\Phi_\sigma=N(0,\sigma^2)\), with \(\Phi_0=\delta_0\), and \(Q_{P_s}^n=Q_s^{\otimes n}\). For \(\sigma>0\), the comparison premise is
\[
\chi^2(Q_s,Q_\star)
\le C\sigma^{-\kappa-1}h^{2\beta+2\kappa+2}
\sum_{j=J}^{\infty}\frac{(\ell^2/\sigma^2)^j}{j!},
\qquad s\in\{+,-\}.
\]
\end{definitionv}

\begin{definitionv}{synth_7}[Chi-squared divergence of observed laws]
For probability measures \(Q\) and \(R\) on the observed-record space \(\{0,1\}\times\mathbb R\times[0,1]\), define
\[
\chi^2(Q,R)=
\begin{cases}
\displaystyle \int\left(\frac{dQ}{dR}-1\right)^2\,dR,
& Q\ll R,\\[4pt]
+\infty,& Q\not\ll R.
\end{cases}
\]
The integral is allowed to equal \(+\infty\). For the witness laws \(P_\pm\) and \(P_\star\), the notation \(\chi^2(P_\pm,P_\star)\) denotes
\[
\chi^2\!\left(\mathcal L_{P_\pm}(X,W,Y),
              \mathcal L_{P_\star}(X,W,Y)\right),
\qquad W=A+\sigma Z,
\]
the divergence between their single-record observed laws.
\end{definitionv}

\begin{definitionv}{synth_8}[Total variation distance]
For probability measures \(\nu_1\) and \(\nu_2\) on a common measurable space \((\Omega,\mathcal F)\), define their total variation distance by
\[
\operatorname{TV}(\nu_1,\nu_2)
=\sup_{B\in\mathcal F}\bigl|\nu_1(B)-\nu_2(B)\bigr|.
\]
\end{definitionv}

\begin{theoremv}{thm:observed-lower}[Explicit observed lower witnesses]*
Fix public class constants \(K=(p_{\min},c_{\mathrm{lo}},c_{\mathrm{hi}})\) and exponents \(\beta,\kappa\) satisfying
\begin{itemize}
\item \textbf{(Class constants.)} \(0<p_{\min}\le1/2\) and \(0<c_{\mathrm{lo}}\le c_{\mathrm{hi}}\).
\item \textbf{(Exponents.)} \(0<\beta\le1\) and \(0\le\kappa\le2\).
\item \textbf{(Design calibration.)}
\[
c_{\mathrm{lo}}\le(\kappa+1)2^\kappa\le c_{\mathrm{hi}}.
\]
\end{itemize}
For every total-variation level \(\tau>0\), there exist constants \(c,C,\epsilon,C_p,c_p>0\) and an integer \(n_0\), chosen uniformly over the public potential-path space, sample size, and error scale, with the following properties.

For every integer \(m\ge4\), the packets of \cref{obj:def:packet-handle} have positive normalization:
\[
K_m(-1)>0\quad(\kappa>0),
\qquad
\widetilde K_m(0)>0\quad(\kappa=0).
\]
Each \(\psi_m\) is even and continuously differentiable on \(\mathbb R\), vanishes outside \([-1,1]\), and satisfies
\[
\psi_m(0)=1,\qquad
|\psi_m(u)|\le C_p,\qquad
|\psi_m'(u)|\le C_pm
\quad(u\in\mathbb R),
\]
\[
\int_{-1}^{1}|u|^\kappa|\psi_m(u)|\,du\le C_pm^{-\kappa-1},
\qquad
\int_{-1}^{1}|u|^\kappa\psi_m(u)^2\,du\le C_pm^{-\kappa-1},
\]
\[
\int_{-1}^{1}|u|^\kappa u^j\psi_m(u)\,du=0
\quad(j\in\mathbb N,\ j<c_pm).
\]

For every public potential-path space \((\mathcal S,\mathcal E)\) of \cref{obj:def:potential-path-space}, every integer \(n\ge n_0\), and every \(\sigma\in[0,1/4]\), there exist structural probability laws \(P_+,P_-,P_\star\) in the model class of \cref{obj:def:model-class}, with public constants \(K\) and parameters \((\beta,\kappa,\sigma)\). All three satisfy \cref{obj:ass:binary-potential-outcomes,obj:ass:rank-uniform,obj:ass:rank-treatment-independence,obj:ass:structural-threshold,obj:ass:iid}.

These laws admit the following common construction. Let \(X\) be the stratum coordinate on \(\mathcal X=\{0,1\}\), \(A\) the latent dose, \(Z\) the standardized Gaussian error, and \(U\) the rank coordinate. Write \(a_0=1/2\) for the evaluation dose and \(t=A-a_0\) for the centered latent dose. Under each law, \(X,A,Z,U\) are independent, with
\[
P(X=0)=P(X=1)=\frac12,\qquad
Z\sim N(0,1),\qquad U\sim\operatorname{Unif}[0,1].
\]
The conditional density \(g_x\) of \(A-a_0\) in either stratum is
\[
g_x(t)=
\begin{cases}
(\kappa+1)2^\kappa|t|^\kappa,& |t|\le1/2,\\
0,& |t|>1/2,
\end{cases}
\qquad x\in\mathcal X.
\]
There are continuous profiles \(\mu_+,\mu_-:[0,1]\to[0,1]\) and a reference profile \(\mu_\star(a)=1/2\). Under \(P_+\), \(P_-\), and \(P_\star\), respectively, the random schedule and realized outcome are
\[
Y(\cdot)=\iota(\mu_+,U),\quad
Y(\cdot)=\iota(\mu_-,U),\quad
Y(\cdot)=\iota(\mu_\star,U),
\qquad
Y=e(Y(\cdot),A).
\]
Here \(e\) and \(\iota\) are the evaluation and threshold-schedule maps of \cref{obj:def:potential-path-space}. In particular, the latent-design marginals agree:
\[
\mathcal L_{P_+}(X,A)=
\mathcal L_{P_-}(X,A)=
\mathcal L_{P_\star}(X,A).
\]

Writing
\[
\theta(P)=\mathbb E_P[Y(a_0)]
\]
for the population causal mean, and using the frontier rate \(\rho_{n,\sigma}\) of \cref{obj:def:frontier-rate} and the observed product experiment \(Q_P^n\) of \cref{obj:def:observed-experiment}, the witnesses satisfy
\[
|\theta(P_+)-\theta(P_-)|\ge c\rho_{n,\sigma},
\qquad
\operatorname{TV}(Q_{P_+}^n,Q_{P_-}^n)\le\tau.
\]

The same profiles have common tuning data \(h,\ell\in\mathbb R\) and \(m,J\in\mathbb N\), where \(h\) is the localization width, \(\ell\) the support radius, \(m\) the packet degree, and \(J\) the cancellation cutoff. Define the direct-regime bandwidth
\[
h_0=n^{-1/(2\beta+\kappa+1)}.
\]
The tuning satisfies
\[
0<h\le1/4,\qquad 0<\ell\le1/4,
\]
\[
\sigma\le h_0\ \Longrightarrow\
h=h_0,\quad \ell=h,\quad J=0,
\]
\[
h_0<\sigma\ \Longrightarrow\
m\ge4,\quad h=\ell/m,\quad
J=\lfloor c_pm\rfloor,\quad J\ge1.
\]
Define the mean perturbation \(\delta\) by
\[
\delta(t)=
\begin{cases}
\epsilon h_0^\beta\bigl(1-(t/h_0)^2\bigr)^2,
  &\sigma\le h_0,\ |t|\le h_0,\\
0,&\sigma\le h_0,\ |t|>h_0,\\
\epsilon h^\beta\psi_m(t/\ell),&h_0<\sigma.
\end{cases}
\]
For every \(a\in[0,1]\),
\[
\mu_\pm(a)=\frac12\pm\delta(a-a_0).
\]
Moreover,
\[
\delta(t)=0\quad(|t|>\ell),
\qquad
\int_{\mathbb R}t^jg_x(t)\delta(t)\,dt=0
\quad(x\in\mathcal X,\ j\in\mathbb N,\ j<J).
\]

For \(\sigma>0\), let \(\chi^2(P_\pm,P_\star)\) denote the chi-squared divergence between the respective one-observation laws induced by \(P_\pm\) and \(P_\star\) in \cref{obj:def:observed-experiment}. Both signs satisfy
\[
\chi^2(P_\pm,P_\star)
\le
C\sigma^{-\kappa-1}h^{2\beta+2\kappa+2}
\sum_{j=J}^{\infty}\frac{(\ell^2/\sigma^2)^j}{j!}.
\]
\end{theoremv}
% CAUSALSMITH-CITED-SCOPE-BEGIN thm:observed-lower
\verificationfootnotetext{\textbf{Formalization scope.} This result uses the cited conclusion \citep{PetrushevXu2005} (Theorem 2.4, equation (2.14)), \citep{https://people.math.sc.edu/pencho/Publications/kpx-06-28-06-web.pdf} (Theorem 2.2, equation (2.4)), \citep{https://dlmf.nist.gov/18.3} (Table 18.3.1, Jacobi row and its caption; orthogonality and squared norm.) and \citep{https://dlmf.nist.gov/5.11.E12} (Equation 5.11.12, restricted to the positive real axis.). The cited source proof is not formalized here: Lean verifies its use and all remaining steps. If that cited conclusion is false, the portions of this result that depend on it are not certified.}
% CAUSALSMITH-CITED-SCOPE-END thm:observed-lower
